\documentclass[11pt,a4paper]{article}
\usepackage[margin=2.2cm]{geometry}
\usepackage{amsmath,amssymb,xcolor,enumitem,booktabs}
\usepackage[normalem]{ulem}
\usepackage[hidelinks]{hyperref}
\setlength{\parskip}{0.6em}\setlength{\parindent}{0pt}
\definecolor{newc}{HTML}{1f6f43}
\newcommand{\new}[1]{\textcolor{newc}{#1}}
\newcommand{\old}[1]{\textcolor{gray}{\sout{#1}}}
\newenvironment{greenblock}{\color{newc}}{}
\newcommand{\dobs}{d_{\mathrm{obs}}}
\title{Proposed change, what a decentralised agent sees and the learning problem}
\date{18 November 2026, version 1}
\begin{document}\sloppy
\maketitle

\section*{How to read this}
Part~1 checks what I told you in chat and corrects two points. Part~2
settles the credit sentence. Part~3 goes through Section V-G and Part~4
through Section V-H of \texttt{3.ProblemFormulation.tex}, showing only the
paragraphs that change. Black text stays, as context. \new{Green} is new
and \old{struck} text goes. Part~5 is the addition to Implementation, Part~6
what else changes, and Part~7 the questions. Nothing is changed in the
thesis yet.

The changes are the three we agreed on.
\begin{enumerate}[nosep]
  \item The decentralised observation matches the code. It also holds the
        role of each vertex in view and where other agents stand.
  \item The credit paragraph leaves the Problem Formulation, and its point
        goes to Implementation at \texttt{eq:credit}.
  \item The learning problem states the game once, drops what repeats
        ``Who decides'', and stops saying that only the observation differs.
\end{enumerate}

\section*{Part 1. What I checked}

\textbf{The credit, confirmed.} In \texttt{training/env.py} the reward of
every agent is replaced by the mean reward of its group
(\texttt{team\_mean}, line 427). The groups are fixed at the start of the
timestep (\texttt{\_groups}, line 470). The centralised controller has one
group, the section-based one a group per zone, and the decentralised
controller is not averaged at all, so each agent keeps its own reward. Every
agent is in exactly one group. So the credits of all agents add up to the
fleet's reward at every timestep, and the fleet average in
\texttt{eq:objective} is the same number under all three controllers.

\textbf{What the credit changes, confirmed.} An agent's choice is judged by
the rewards of its group. With the fleet mean, a choice that holds up
another agent lowers the number it is judged by. With its own reward, it
does not. So in principle centralised training follows the objective's own
direction of improvement, up to PPO's approximations, and decentralised
training leaves out the cost an agent puts on others.

\textbf{Correction 1, the reason for the credit.} I said the credit was set
this way because it is the effect the comparison should show, citing [A9]
and [A32]. That is not their reason. [A9] says credit follows the agents
that choose together ``because they act as one decision-maker'', and that
the commons question is no longer part of the study. The phrase ``the effect
the comparison is meant to show'' is from the critic paragraph in
Implementation, which is about the critic, not the credit.

\textbf{Correction 2, the decentralised agent is not blind to crowding.} I
said a decentralised agent ``never pays for blocking'' another. It does pay
a little. Its own reward has a crowding term, $-n_{\mathrm{cng}}\,\delta_i(t)$
(\texttt{reward\_function.py}, line 8), so a crowd around it costs it
something. What it does not pay for is the time the agents it holds up
lose.

\textbf{Other claims from the margin comments.}
\begin{itemize}[nosep]
  \item The observation in the code holds, for each vertex in view, its role,
        the distance label, whether another agent stands there and whether an
        earlier agent claimed it (\texttt{observation.py}, lines 36 to 50).
        The claim flag is always empty for a decentralised agent. Confirmed.
  \item The dock of a put-away and of a shipment is drawn at random in every
        fleet config (\texttt{dock\_choice: random}). The Problem
        Formulation already says so in ``Orders and tasks'', so it now
        contradicts itself, in V-H and in Scope (``the only randomness is
        which orders arrive''). Confirmed.
  \item One priority order, drawn once per episode, is used by the
        centralised turns, conflict resolution and assignment
        (\texttt{multi\_agent\_env.py}, lines 517, 614, 1009). So the
        ``priority order'' in V-H is the one of ``Who decides''.
  \item Figure~9. The script says $a_1$ ``is also heading for W1'' on
        purpose. The text never says why, so the comment stands.
  \item Figure~10 uses a field of view of one step (\texttt{crowding.py}).
        Confirmed.
  \item Two of my margin comments were too strong. Every generated warehouse
        sets each edge cost to 1, so ``fewest steps'' and ``travel distance''
        are the same number, and only the wording differs. The code also
        checks every window against a fixed largest size, so ``a small field
        of view keeps the input the same size'' is close to true. ``Never
        larger than a fixed size'' would be exact.
\end{itemize}

\section*{Part 2. Keep the credit sentence or not}

\textbf{Drop it, with nothing left in its place.} There are three reasons.
\begin{enumerate}[nosep]
  \item How a choice is credited is part of how the policy is trained. That
        is method, and the formula it needs, \texttt{eq:credit}, is in
        Implementation. In the Problem Formulation the claim stands without
        the formula, and the reader cannot check it.
  \item The short form I offered, ``a decentralised agent sees only its own
        view, so it cannot see the delay it causes others'', is not true. The
        agent sees the agents near it, and with change~1 the text says so. It
        also pays for crowding through its reward.
  \item The problem itself does not depend on the credit. The objective is
        the same fleet average for every controller, so the Problem
        Formulation is complete without it.
\end{enumerate}
Its point does matter for the results, so it moves to Implementation
(Part~5), next to the formula it follows from.

\section*{Part 3. Section V-G, what a decentralised agent sees}

\subsection*{Opening}
A decentralised agent does not see the whole warehouse. At every timestep it
sees \old{three}\new{four} things\old{:}\new{.} \new{These are} what is near
it, \new{where the other agents near it stand,} which way its target is, and
how crowded it is around it. Together these form its observation $o_i(t)$. It
may also hear messages from nearby agents.

\subsection*{The field of view}
The agent sees every vertex it can reach within a travel cost of $\dobs$, and
the edges between them. So it knows which moves \old{are open}\new{it can make}
nearby, not only which vertices exist. \new{It also sees the role of each of
these vertices, whether it is a shelf, a workstation, a dock, an endpoint or
an aisle position.} It sees nothing further away. Formally, the \emph{field of
view} of agent $a_i$ is the subgraph of $G$ on those vertices,
\[
  V^{(\dobs)}_i(t) = \{ v \in V : d_G(\ell_i(t), v) \le \dobs \}, \qquad
  \mathcal{G}^{(\dobs)}_i(t) = G\bigl[V^{(\dobs)}_i(t)\bigr]. \tag{18}
\]

\begin{greenblock}
Knowing where the aisles go is not enough to avoid other agents. So the agent
also sees which vertices in its field of view hold another agent. It sees
that a vertex is taken, not which agent stands there or where that agent is
heading. Formally, the \emph{occupied vertices} in view are
\[
  L_i(t) = \{ \ell_j(t) : a_j \in \mathcal{A}\setminus\{a_i\} \}
           \cap V^{(\dobs)}_i(t) .
\]
\end{greenblock}

Figure~9 shows it for agent $a_1$ with $\dobs=2$, in the warehouse of the tea
task $\tau_1$ at $t=12$, after $a_2$ has picked up $\tau_1$ at $A$. \new{Agent
$a_1$ carries another task to the same workstation $W_1$.} The rest of this
section uses this agent $a_1$ as its example. \new{Two of the five vertices
around it are occupied, those of $a_2$ and $a_3$.} \old{Communication uses the
same range, so the field of view is also how far $a_1$ can hear.}

\subsection*{Direction to the target, the example paragraph}
A smaller $\eta_i(v,t)$ means $v$ is closer to the target. In Figure~9, agent
$a_1$ can step left to $v$ or right to $w$, with $\eta_1(v,t)=3$ and
$\eta_1(w,t)=5$. $v$ leaves less distance to cover, but $a_2$ stands on the
short route just past $v$\new{, and $a_1$ sees that this vertex is occupied}.
A trained policy may therefore still send $a_1$ to $w$, or make it wait, if it
expects the route through $v$ to be congested.

\subsection*{How crowded the neighbourhood is}
No text changes, apart from one sentence after the formula.

\new{The crowding feature is the share of the field of view that $L_i(t)$
covers. The agent could count it itself, and it is given as one number
because it is also used outside the observation.}

(See Question~3. If you prefer, this sentence is left out.)

\subsection*{Putting the pieces together}
The \old{three}\new{four} pieces together form the \emph{decentralised
observation} of agent $a_i$,
\[
  o_i(t) = \bigl( \mathcal{G}^{(\dobs)}_i(t),\; \new{L_i(t),}\; \eta_i(\cdot,
  t),\; \delta_i(t) \bigr). \tag{23}
\]

The messages subsection says ``$o_i(t)$ keeps its three pieces'', which
becomes ``\new{four}''.

\section*{Part 4. Section V-H, the learning problem}

\subsection*{The cycle}
Every timestep follows the same cycle. Each agent looks at what its
controller's window shows, all agents act at the same time, the warehouse
moves to its next state, and each agent receives a number saying how good
that step was. Except under the centralised controller, an agent neither sees
the whole warehouse nor knows what every other agent will do. A partially
observable stochastic game is the standard model of this situation
(Hansen et al.), so the learning problem is written as one.

\subsection*{The game, stated once}
\old{Every part of the game is something the reader has already met. The
agents are the fleet. The state $s_t$ is the one of Section V-F, Who decides.
It also holds the priority order that conflict resolution uses. The actions
are wait and move. The transition rule turns a state and the agents' actions
into the next state. The observation is what each agent sees. The reward
scores one step for one agent, and the discount sets how much a later reward
counts against one now.}

\new{The state $s_t$ is the one of ``Who decides'', together with the
priority order drawn there, which conflict resolution and assignment also
use.} Formally, the learning problem is
\[
  \mathcal{M} = \bigl( \mathcal{A},\, \mathcal{S},\, \{\mathcal{U}_i\},\, P,\,
  \{\mathcal{O}_i\},\, \{O_i\},\, \{R_i\},\, \gamma \bigr), \tag{24}
\]
\old{with the agents $\mathcal{A}$, the state space $\mathcal{S}$, the action
sets $\mathcal{U}_i$, the transition rule $P$, the observation spaces
$\mathcal{O}_i$ and functions $O_i$, the rewards $R_i$ and the discount
$\gamma$. The reward of agent $a_i$, $R_i : \mathcal{S} \times \mathcal{U}_1
\times \dots \times \mathcal{U}_m \times \mathcal{S} \to \mathbb{R}$, takes the
state before the step, the joint action and the state after it.}
\new{where Table~V gives each element and where it was defined.}

Table~V stays as it is.

\subsection*{The action list}
A stochastic game gives every agent one fixed list of actions, but the legal
moves depend on where the agent stands. So the list is as long as the busiest
vertex needs, one wait and one move per outgoing edge. Move number $k$ takes
the $k$-th outgoing edge of the vertex the agent stands on, in vertex-number
order. Moves that lead nowhere are switched off before the policy chooses
\old{(Section VI-C)}\new{, so the ones left are the legal actions
$\mathcal{U}(\ell_i(t))$ of ``Who decides''}. If the busiest junction has four
exits, every agent has five actions, and in a one-way aisle two of them are
switched on. Formally, with $d_{\max}$ the largest number of outgoing edges at
any vertex, the action list is $\mathcal{U}_i = \{\mathrm{wait},
\mathrm{move}_1, \dots, \mathrm{move}_{d_{\max}}\}$.

\subsection*{The policy}
The policy and the transition \old{kernel}\new{rule} answer two different
questions. The policy answers what each agent tries to do. From what the agent
observes it gives a probability to each legal action, for example
\old{move north}\new{$\mathrm{move}_1$} with probability $0.7$ and wait with
probability $0.3$, and one action is drawn. The drawn actions of all agents,
side by side, form the joint action $u_t$, the whole fleet's decision for this
timestep. \old{Under the centralised controller the agents draw their actions
one after another, each seeing the earlier draws, and the result is the same
kind of joint action. The section-based controller does the same inside each
zone.}

\subsection*{The transition rule}
The transition \old{kernel}\new{rule} answers what then happens to the
warehouse. It takes the current state and the joint action and applies
everything that is not a choice of any agent. Conflicts are resolved, the
shared assignment rule pairs free agents with waiting tasks, the storage
update is applied, and newly arrived orders are added. The result is the next
state $s_{t+1}$. \old{The kernel is therefore not a state but the rule that
produces one. Its only randomness is the arrival of orders,}\new{It is random
in two places only, which orders arrive and which dock a put-away or a
shipment is drawn to,} so the same state and joint action always move the
agents in the same way. Formally, $P(s' \mid s_t, u_t)$ is the probability
that the next state is $s'$. \old{The assignment rule and conflict resolution
are the same for every controller, so they are part of $P$, not of any
policy.}

\subsection*{The observation function}
\old{Only the observation changes between the controllers. Everything else in
the game is the same.}\new{The controllers differ in two things, what an agent
observes and which agents choose together. Which agents choose together is
part of the policy, as set out in ``Who decides''. What an agent observes is
the observation function.} A centralised agent sees the whole state, a
section-based agent sees its zone, and a decentralised agent sees its own
view. Formally, the \emph{observation function} is
\[
  O_i(s_t) =
  \begin{cases}
    (s_t,\, i) & \text{centralised},\\
    (s^{(q)}_t,\, i) \text{ with } \ell_i(t) \in Y_q & \text{section-based},\\
    o_i(t) & \text{decentralised}.
  \end{cases} \tag{25}
\]
\old{It is deterministic, so the same state always gives the same
observation.}

The three views fit inside each other. The centralised controller sees
everything the section-based one sees, and that one sees everything a
decentralised agent sees. \old{Messages and the moves of earlier agents are
not part of what an agent sees, so they do not change this.}\new{The moves of
earlier agents and the messages are inputs to the policy and not part of
$O_i$, so they are left out of this comparison.}

\subsection*{The objective}
All agents are alike, work toward the same goal and share one set of
parameters $\theta$. \new{For the decentralised controller these are the
weights every agent uses. For the centralised and section-based controllers
they are the weights of the network over the whole graph or over a zone.} The
learning objective is therefore stated for the whole team. It is the average
reward the fleet earns per timestep, added up over the horizon, with rewards
further in the future counting less. The discount factor $\gamma \in [0,1)$
sets how much less,
\[
  \max_{\theta} J(\theta) = \max_{\theta} \mathbb{E}_{\pi_\theta}
  \Bigl[ \sum_{t=0}^{T-1} \gamma^{t}\, \frac{1}{m} \sum_{i=1}^{m} R_i(t)
  \Bigr]. \tag{26}
\]
Here $\pi_\theta$ stands for the policy being trained,
$\pi^{\mathrm{cen}}_\theta$, $\pi^{\mathrm{sec}}_\theta$ or
$\pi^{\mathrm{dec}}_\theta$, by the controller. $R_i(t)$ is short for
$R_i(s_t, u_t, s_{t+1})$, the reward agent $a_i$ earns on the step from $s_t$
to $s_{t+1}$.

\old{The same objective is used for all three learned controllers. Only what
$\theta$ stands for changes. For the decentralised controller it is the
weights shared by all agents. For the centralised and section-based
controllers it is the weights of the network over the whole graph or over a
zone, as in (14) and (16).}

\old{Sharing one objective does not mean the three controllers learn the same
thing. An action is judged by the average reward of the agents that chose
together with it. So a decentralised agent is not charged for the delay it
causes others, but a centralised agent is.}

The reward $R_i$ itself is left open here. It is only a tool for getting a
policy that does well on the system-level measures of Section V-J, so it can
be redesigned without changing the problem. It is given in Section VI-F.

\section*{Part 5. Implementation, after \texttt{eq:credit}}

The paragraph that ends with \texttt{eq:credit} gets two new paragraphs after
the formula.

\begin{greenblock}
The groups split the fleet, so the credits of all agents add up to the
fleet's reward at every timestep. Every controller therefore works toward the
same objective, the fleet average of eq.~(26). The credit only changes which
rewards an agent's choice is judged by.

A centralised agent's choice is judged by the reward of the whole fleet, so
it pays when it holds up any other agent. A section-based agent pays only
when the agent it holds up is in its own zone. A decentralised agent is
judged by its own reward. It pays for the crowd around it through the
crowding term, but not for the time it costs the agents it holds up. Credit
follows the agents that choose together because they act as one
decision-maker.
\end{greenblock}

\section*{Part 6. What changes elsewhere}
\begin{itemize}[nosep]
  \item \texttt{A.Notation.tex} and \texttt{2.Notation.tex}, a new row for
        $L_i(t)$, ``vertices in agent $i$'s field of view where another agent
        stands''. The letter $L$ is free.
  \item Scope and assumptions, ``Deterministic execution'': ``the only
        randomness is which orders arrive'' becomes ``the only randomness is
        which orders arrive and which dock a put-away or shipment uses''.
  \item \texttt{GAPS.md}, a new entry [A107] with the three changes, the two
        tags \texttt{\% [A107]} in the chapters, and a note that it revises
        the randomness sentence of [A12]/[A32].
  \item Figure~9 is not redrawn. It already shows where the agents stand.
\end{itemize}

\section*{Part 7. Questions}
\begin{enumerate}
  \item Is $L_i(t)$ worth a symbol, or should the occupied vertices stay in
        words and the tuple read ``the field of view with its occupied
        vertices''? I added it because the tuple in eq.~(23) needs something to
        list.
  \item The code also gives a decentralised agent the role of each vertex. I
        put that in words inside the field of view, since the roles are
        part of $G$ already. Is that enough, or should the role show in the
        tuple too?
  \item With $L_i(t)$ in the observation, $\delta_i(t)$ can be counted from it.
        Keep the sentence in Part~3 that says so, or leave it out?
  \item The other margin comments are not in this proposal. These are the
        messages placed before ``Putting the pieces together'', the idle
        target paragraph's place, the forward references to Section VI, the
        names in Figure~10, and the reasons for $\dobs=2$ and for hearing
        along the aisles. Take them in a second round, or add them here?
\end{enumerate}

\end{document}
